%%%%%%%%%%%%%%%%%%%%%%%%%%%%%%%%%%%%%%%%%%%%%%%%%%%%%%%%%%%%
% 11.4 NUMERICAL ORDINARY DIFFERENTIAL EQUATIONS
% The Composition of Advanced Mathematics
%%%%%%%%%%%%%%%%%%%%%%%%%%%%%%%%%%%%%%%%%%%%%%%%%%%%%%%%%%%%

\section{Numerical Ordinary Differential Equations}
\label{sec:numerical-odes}

\prerequisite{
Ordinary differential equations, initial-value problems, Taylor series,
numerical analysis, numerical linear algebra, nonlinear equation solving,
and basic programming concepts.
}

\learningobjectives{
By the end of this section, the reader should be able to:
\begin{itemize}
    \item formulate an ordinary differential-equation initial-value problem;
    \item construct a discrete time grid;
    \item derive the forward Euler method;
    \item interpret Euler's method geometrically;
    \item distinguish local and global truncation error;
    \item understand consistency and convergence of ODE solvers;
    \item understand zero-stability conceptually;
    \item derive backward Euler;
    \item distinguish explicit and implicit time-stepping methods;
    \item recognize the trapezoidal method;
    \item derive second-order Runge--Kutta methods;
    \item understand midpoint and improved Euler methods;
    \item apply the classical fourth-order Runge--Kutta method;
    \item interpret Runge--Kutta stages;
    \item understand Butcher-tableau notation;
    \item distinguish one-step and multistep methods;
    \item understand Adams--Bashforth methods;
    \item understand Adams--Moulton methods;
    \item recognize predictor--corrector methods;
    \item understand backward differentiation formulas conceptually;
    \item define absolute stability;
    \item analyze methods using the linear test equation;
    \item determine stability functions;
    \item understand stability regions;
    \item distinguish A-stability and L-stability conceptually;
    \item recognize stiff differential equations;
    \item explain why explicit methods struggle with stiffness;
    \item understand adaptive time stepping;
    \item recognize embedded Runge--Kutta pairs;
    \item understand local error estimation;
    \item choose time steps using error estimates;
    \item understand relative and absolute ODE tolerances;
    \item solve systems of ODEs numerically;
    \item convert higher-order ODEs to first-order systems;
    \item solve implicit method equations using Newton iteration;
    \item understand Jacobians in stiff solvers;
    \item recognize event detection;
    \item distinguish interpolation from dense output;
    \item understand numerical conservation issues;
    \item recognize symplectic methods conceptually;
    \item recognize positivity and invariant preservation issues;
    \item understand shooting methods for boundary-value problems;
    \item recognize finite-difference methods for ODE boundary-value problems;
    \item understand method-of-lines connections to PDEs;
    \item perform convergence studies;
    \item connect numerical ODEs with simulation, control, mechanics,
          biology, chemistry, economics, and scientific computing.
\end{itemize}
}

%%%%%%%%%%%%%%%%%%%%%%%%%%%%%%%%%%%%%%%%%%%%%%%%%%%%%%%%%%%%
\subsection{Initial-Value Problems}
\label{subsec:numerical-ode-ivp}
%%%%%%%%%%%%%%%%%%%%%%%%%%%%%%%%%%%%%%%%%%%%%%%%%%%%%%%%%%%%

A first-order ordinary differential-equation initial-value problem has the
form

\[
\boxed{
y'(t)
=
f(t,y(t)),
}
\]

with

\[
\boxed{
y(t_0)
=
y_0.
}
\]

The goal is to approximate the solution at a collection of time points.

%%%%%%%%%%%%%%%%%%%%%%%%%%%%%%%%%%%%%%%%%%%%%%%%%%%%%%%%%%%%
\subsection{Time Grid}
\label{subsec:ode-time-grid}
%%%%%%%%%%%%%%%%%%%%%%%%%%%%%%%%%%%%%%%%%%%%%%%%%%%%%%%%%%%%

Choose grid points

\[
\boxed{
t_n
=
t_0+nh,
}
\]

where

\[
h>0
\]

is the time step.

A numerical method produces approximations

\[
\boxed{
y_n
\approx
y(t_n).
}
\]

For variable-step methods, the step size may instead depend on \(n\):

\[
\boxed{
h_n
=
t_{n+1}-t_n.
}
\]

%%%%%%%%%%%%%%%%%%%%%%%%%%%%%%%%%%%%%%%%%%%%%%%%%%%%%%%%%%%%
\subsection{Forward Euler Method}
\label{subsec:forward-euler}
%%%%%%%%%%%%%%%%%%%%%%%%%%%%%%%%%%%%%%%%%%%%%%%%%%%%%%%%%%%%

From Taylor expansion,

\[
y(t+h)
=
y(t)
+
hy'(t)
+
O(h^2).
\]

Since

\[
y'(t)
=
f(t,y(t)),
\]

we obtain

\[
y(t+h)
=
y(t)
+
hf(t,y(t))
+
O(h^2).
\]

Ignoring the higher-order term gives the forward Euler method:

\[
\boxed{
y_{n+1}
=
y_n
+
h
f(t_n,y_n).
}
\]

%%%%%%%%%%%%%%%%%%%%%%%%%%%%%%%%%%%%%%%%%%%%%%%%%%%%%%%%%%%%
\subsection{Geometric Interpretation of Euler's Method}
\label{subsec:euler-geometric}
%%%%%%%%%%%%%%%%%%%%%%%%%%%%%%%%%%%%%%%%%%%%%%%%%%%%%%%%%%%%

At point

\[
(t_n,y_n),
\]

the differential equation provides slope

\[
\boxed{
f(t_n,y_n).
}
\]

Euler's method follows the tangent line for time \(h\):

\[
\boxed{
y_{n+1}
=
y_n
+
h
\times
\text{slope}.
}
\]

Thus Euler's method repeatedly approximates the solution curve by short
tangent-line segments.

%%%%%%%%%%%%%%%%%%%%%%%%%%%%%%%%%%%%%%%%%%%%%%%%%%%%%%%%%%%%
\subsection{Worked Example: Forward Euler}
\label{subsec:worked-forward-euler}
%%%%%%%%%%%%%%%%%%%%%%%%%%%%%%%%%%%%%%%%%%%%%%%%%%%%%%%%%%%%

\begin{workedexamplebox}{Approximating an Exponential Solution}

Consider

\[
y'=y,
\qquad
y(0)=1.
\]

The exact solution is

\[
y(t)=e^t.
\]

Take

\[
h=0.1.
\]

Euler's method gives

\[
y_{n+1}
=
y_n+0.1y_n
=
1.1y_n.
\]

Starting from

\[
y_0=1,
\]

we obtain

\[
y_1=1.1,
\]

\[
y_2=1.21,
\]

and

\[
y_3=1.331.
\]

The exact value at

\[
t=0.3
\]

is

\[
e^{0.3}
\approx1.34986.
\]

\begin{answerbox}
\[
\boxed{
y_3=1.331
}
\]

compared with

\[
\boxed{
y(0.3)\approx1.34986.
}
\]
\end{answerbox}

\end{workedexamplebox}

%%%%%%%%%%%%%%%%%%%%%%%%%%%%%%%%%%%%%%%%%%%%%%%%%%%%%%%%%%%%
\subsection{Local Truncation Error}
\label{subsec:ode-local-truncation-error}
%%%%%%%%%%%%%%%%%%%%%%%%%%%%%%%%%%%%%%%%%%%%%%%%%%%%%%%%%%%%

The local truncation error measures the error introduced in one step when
the exact solution is used as the starting value.

For forward Euler,

\[
y(t_{n+1})
=
y(t_n)
+
hf(t_n,y(t_n))
+
O(h^2).
\]

Therefore the one-step defect is

\[
\boxed{
O(h^2).
}
\]

Depending on convention, the local truncation error may be divided by
\(h\), giving an \(O(h)\) quantity.

The convention should always be stated clearly.

%%%%%%%%%%%%%%%%%%%%%%%%%%%%%%%%%%%%%%%%%%%%%%%%%%%%%%%%%%%%
\subsection{Global Error}
\label{subsec:ode-global-error}
%%%%%%%%%%%%%%%%%%%%%%%%%%%%%%%%%%%%%%%%%%%%%%%%%%%%%%%%%%%%

The global error is

\[
\boxed{
e_n
=
y(t_n)-y_n.
}
\]

Forward Euler accumulates local errors across approximately

\[
O(1/h)
\]

steps over a fixed time interval.

Under standard smoothness and stability assumptions,

\[
\boxed{
e_n
=
O(h).
}
\]

Thus forward Euler is a first-order method.

%%%%%%%%%%%%%%%%%%%%%%%%%%%%%%%%%%%%%%%%%%%%%%%%%%%%%%%%%%%%
\subsection{Order of an ODE Method}
\label{subsec:ode-method-order}
%%%%%%%%%%%%%%%%%%%%%%%%%%%%%%%%%%%%%%%%%%%%%%%%%%%%%%%%%%%%

A numerical ODE method has global order \(p\) if

\[
\boxed{
\max_{0\leq n\leq N}
|y(t_n)-y_n|
=
O(h^p)
}
\]

over a fixed interval as

\[
h\to0.
\]

If the step size is halved, the global error should decrease by
approximately

\[
\boxed{
2^p
}
\]

in the asymptotic regime.

%%%%%%%%%%%%%%%%%%%%%%%%%%%%%%%%%%%%%%%%%%%%%%%%%%%%%%%%%%%%
\subsection{Consistency}
\label{subsec:ode-consistency}
%%%%%%%%%%%%%%%%%%%%%%%%%%%%%%%%%%%%%%%%%%%%%%%%%%%%%%%%%%%%

An ODE discretization is consistent if the discrete approximation tends
toward the differential equation as

\[
h\to0.
\]

For forward Euler,

\[
\frac{
y_{n+1}-y_n
}{
h
}
=
f(t_n,y_n).
\]

As \(h\to0\), the difference quotient approaches

\[
y'(t).
\]

%%%%%%%%%%%%%%%%%%%%%%%%%%%%%%%%%%%%%%%%%%%%%%%%%%%%%%%%%%%%
\subsection{Stability and Convergence}
\label{subsec:ode-stability-convergence}
%%%%%%%%%%%%%%%%%%%%%%%%%%%%%%%%%%%%%%%%%%%%%%%%%%%%%%%%%%%%

Consistency alone is insufficient.

Numerical perturbations and local truncation errors must not grow
uncontrollably.

For many classes of ODE methods, convergence depends on both:

\[
\boxed{
\text{consistency}
}
\]

and

\[
\boxed{
\text{stability}.
}
\]

For linear multistep methods, the relevant stability notion includes
zero-stability.

%%%%%%%%%%%%%%%%%%%%%%%%%%%%%%%%%%%%%%%%%%%%%%%%%%%%%%%%%%%%
\subsection{Backward Euler Method}
\label{subsec:backward-euler}
%%%%%%%%%%%%%%%%%%%%%%%%%%%%%%%%%%%%%%%%%%%%%%%%%%%%%%%%%%%%

Backward Euler evaluates the slope at the new time:

\[
\boxed{
y_{n+1}
=
y_n
+
h
f(t_{n+1},y_{n+1}).
}
\]

Because \(y_{n+1}\) appears inside \(f\), this is generally an implicit
equation.

%%%%%%%%%%%%%%%%%%%%%%%%%%%%%%%%%%%%%%%%%%%%%%%%%%%%%%%%%%%%
\subsection{Explicit and Implicit Methods}
\label{subsec:explicit-implicit-ode-methods}
%%%%%%%%%%%%%%%%%%%%%%%%%%%%%%%%%%%%%%%%%%%%%%%%%%%%%%%%%%%%

An explicit method computes the new state directly from known quantities.

For example,

\[
\boxed{
y_{n+1}
=
y_n
+
hf(t_n,y_n).
}
\]

An implicit method requires solving an equation involving the unknown new
state.

For example,

\[
\boxed{
y_{n+1}
-
y_n
-
hf(t_{n+1},y_{n+1})
=
0.
}
\]

%%%%%%%%%%%%%%%%%%%%%%%%%%%%%%%%%%%%%%%%%%%%%%%%%%%%%%%%%%%%
\subsection{Worked Example: Backward Euler}
\label{subsec:worked-backward-euler}
%%%%%%%%%%%%%%%%%%%%%%%%%%%%%%%%%%%%%%%%%%%%%%%%%%%%%%%%%%%%

\begin{workedexamplebox}{Backward Euler for Exponential Decay}

Consider

\[
y'=-5y,
\qquad
y(0)=1.
\]

Backward Euler gives

\[
y_{n+1}
=
y_n
-
5hy_{n+1}.
\]

Therefore,

\[
(1+5h)y_{n+1}
=
y_n.
\]

Hence,

\[
\boxed{
y_{n+1}
=
\frac{
y_n
}{
1+5h
}.
}
\]

For

\[
h=0.1,
\]

\[
y_1
=
\frac1{1.5}
=
\frac23.
\]

\begin{answerbox}
\[
\boxed{
y_1
\approx0.6667.
}
\]
\end{answerbox}

\end{workedexamplebox}

%%%%%%%%%%%%%%%%%%%%%%%%%%%%%%%%%%%%%%%%%%%%%%%%%%%%%%%%%%%%
\subsection{Implicit Solve for Nonlinear ODEs}
\label{subsec:implicit-nonlinear-ode-solve}
%%%%%%%%%%%%%%%%%%%%%%%%%%%%%%%%%%%%%%%%%%%%%%%%%%%%%%%%%%%%

For nonlinear \(f\), backward Euler requires solving

\[
\boxed{
G(z)
=
z-y_n-hf(t_{n+1},z)
=
0.
}
\]

Newton's method gives

\[
\boxed{
z^{(m+1)}
=
z^{(m)}
-
[
G'(z^{(m)})
]^{-1}
G(z^{(m)}).
}
\]

Thus implicit ODE integration may require nonlinear optimization or
root-finding technology at every time step.

%%%%%%%%%%%%%%%%%%%%%%%%%%%%%%%%%%%%%%%%%%%%%%%%%%%%%%%%%%%%
\subsection{Trapezoidal Method}
\label{subsec:ode-trapezoidal-method}
%%%%%%%%%%%%%%%%%%%%%%%%%%%%%%%%%%%%%%%%%%%%%%%%%%%%%%%%%%%%

Average the slopes at the beginning and end of the interval:

\[
\boxed{
y_{n+1}
=
y_n
+
\frac{h}{2}
\left[
f(t_n,y_n)
+
f(t_{n+1},y_{n+1})
\right].
}
\]

This is an implicit second-order method.

It is sometimes called the Crank--Nicolson method in closely related
time-discretization contexts.

%%%%%%%%%%%%%%%%%%%%%%%%%%%%%%%%%%%%%%%%%%%%%%%%%%%%%%%%%%%%
\subsection{Improved Euler Method}
\label{subsec:improved-euler}
%%%%%%%%%%%%%%%%%%%%%%%%%%%%%%%%%%%%%%%%%%%%%%%%%%%%%%%%%%%%

An explicit predictor is

\[
\boxed{
\widetilde{y}_{n+1}
=
y_n
+
hf(t_n,y_n).
}
\]

Then average the initial and predicted final slopes:

\[
\boxed{
y_{n+1}
=
y_n
+
\frac{h}{2}
\left[
f(t_n,y_n)
+
f(t_{n+1},\widetilde{y}_{n+1})
\right].
}
\]

This method is often called Heun's method or the explicit trapezoidal
method.

%%%%%%%%%%%%%%%%%%%%%%%%%%%%%%%%%%%%%%%%%%%%%%%%%%%%%%%%%%%%
\subsection{Explicit Midpoint Method}
\label{subsec:explicit-midpoint}
%%%%%%%%%%%%%%%%%%%%%%%%%%%%%%%%%%%%%%%%%%%%%%%%%%%%%%%%%%%%

First estimate the midpoint:

\[
\boxed{
y_{n+1/2}
=
y_n
+
\frac{h}{2}
f(t_n,y_n).
}
\]

Then use the midpoint slope:

\[
\boxed{
y_{n+1}
=
y_n
+
h
f
\left(
t_n+\frac{h}{2},
y_{n+1/2}
\right).
}
\]

This is a second-order Runge--Kutta method.

%%%%%%%%%%%%%%%%%%%%%%%%%%%%%%%%%%%%%%%%%%%%%%%%%%%%%%%%%%%%
\subsection{Runge--Kutta Methods}
\label{subsec:runge-kutta-methods}
%%%%%%%%%%%%%%%%%%%%%%%%%%%%%%%%%%%%%%%%%%%%%%%%%%%%%%%%%%%%

Runge--Kutta methods approximate the effect of several slope samples within
each time step.

A general explicit \(s\)-stage method has stages

\[
\boxed{
k_i
=
f
\left(
t_n+c_ih,
y_n
+
h
\sum_{j=1}^{i-1}
a_{ij}k_j
\right),
}
\]

and update

\[
\boxed{
y_{n+1}
=
y_n
+
h
\sum_{i=1}^{s}
b_ik_i.
}
\]

%%%%%%%%%%%%%%%%%%%%%%%%%%%%%%%%%%%%%%%%%%%%%%%%%%%%%%%%%%%%
\subsection{Classical Fourth-Order Runge--Kutta Method}
\label{subsec:rk4}
%%%%%%%%%%%%%%%%%%%%%%%%%%%%%%%%%%%%%%%%%%%%%%%%%%%%%%%%%%%%

The classical fourth-order Runge--Kutta method, abbreviated RK4, uses

\[
\boxed{
k_1
=
f(t_n,y_n),
}
\]

\[
\boxed{
k_2
=
f
\left(
t_n+\frac{h}{2},
y_n+\frac{h}{2}k_1
\right),
}
\]

\[
\boxed{
k_3
=
f
\left(
t_n+\frac{h}{2},
y_n+\frac{h}{2}k_2
\right),
}
\]

and

\[
\boxed{
k_4
=
f
(
t_n+h,
y_n+hk_3
).
}
\]

Then

\[
\boxed{
y_{n+1}
=
y_n
+
\frac{h}{6}
(
k_1+2k_2+2k_3+k_4
).
}
\]

%%%%%%%%%%%%%%%%%%%%%%%%%%%%%%%%%%%%%%%%%%%%%%%%%%%%%%%%%%%%
\subsection{Interpretation of RK4}
\label{subsec:rk4-interpretation}
%%%%%%%%%%%%%%%%%%%%%%%%%%%%%%%%%%%%%%%%%%%%%%%%%%%%%%%%%%%%

RK4 samples slopes at:

\[
\boxed{
\text{the beginning},
}
\]

\[
\boxed{
\text{two midpoint estimates},
}
\]

and

\[
\boxed{
\text{the endpoint}.
}
\]

The weighted average

\[
\frac16
(
k_1+2k_2+2k_3+k_4
)
\]

approximates the average derivative over the step.

%%%%%%%%%%%%%%%%%%%%%%%%%%%%%%%%%%%%%%%%%%%%%%%%%%%%%%%%%%%%
\subsection{Worked Example: RK4}
\label{subsec:worked-rk4}
%%%%%%%%%%%%%%%%%%%%%%%%%%%%%%%%%%%%%%%%%%%%%%%%%%%%%%%%%%%%

\begin{workedexamplebox}{One RK4 Step}

Consider

\[
y'=y,
\qquad
y(0)=1,
\]

with

\[
h=0.1.
\]

Then

\[
k_1=1.
\]

Next,

\[
k_2
=
1+0.05(1)
=
1.05.
\]

Then,

\[
k_3
=
1+0.05(1.05)
=
1.0525.
\]

Finally,

\[
k_4
=
1+0.1(1.0525)
=
1.10525.
\]

Thus,

\[
y_1
=
1
+
\frac{0.1}{6}
\left(
1
+
2(1.05)
+
2(1.0525)
+
1.10525
\right).
\]

Therefore,

\[
y_1
\approx1.1051708.
\]

The exact value is

\[
e^{0.1}
\approx1.1051702.
\]

\begin{answerbox}
\[
\boxed{
y_1\approx1.1051708.
}
\]
\end{answerbox}

\end{workedexamplebox}

%%%%%%%%%%%%%%%%%%%%%%%%%%%%%%%%%%%%%%%%%%%%%%%%%%%%%%%%%%%%
\subsection{Butcher Tableaux}
\label{subsec:butcher-tableaux}
%%%%%%%%%%%%%%%%%%%%%%%%%%%%%%%%%%%%%%%%%%%%%%%%%%%%%%%%%%%%

A Runge--Kutta method can be represented by a Butcher tableau:

\[
\boxed{
\begin{array}{c|cccc}
c_1 & a_{11} & a_{12} & \cdots & a_{1s}\\
c_2 & a_{21} & a_{22} & \cdots & a_{2s}\\
\vdots & \vdots & \vdots & & \vdots\\
c_s & a_{s1} & a_{s2} & \cdots & a_{ss}\\
\hline
& b_1 & b_2 & \cdots & b_s
\end{array}
}
\]

For an explicit Runge--Kutta method,

\[
\boxed{
a_{ij}=0
\qquad
\text{for}
\qquad
j\geq i.
}
\]

%%%%%%%%%%%%%%%%%%%%%%%%%%%%%%%%%%%%%%%%%%%%%%%%%%%%%%%%%%%%
\subsection{Euler Butcher Tableau}
\label{subsec:euler-butcher-tableau}
%%%%%%%%%%%%%%%%%%%%%%%%%%%%%%%%%%%%%%%%%%%%%%%%%%%%%%%%%%%%

Forward Euler has tableau

\[
\boxed{
\begin{array}{c|c}
0&0\\
\hline
&1
\end{array}
}
\]

because it uses one stage:

\[
k_1=f(t_n,y_n).
\]

%%%%%%%%%%%%%%%%%%%%%%%%%%%%%%%%%%%%%%%%%%%%%%%%%%%%%%%%%%%%
\subsection{Classical RK4 Tableau}
\label{subsec:rk4-tableau}
%%%%%%%%%%%%%%%%%%%%%%%%%%%%%%%%%%%%%%%%%%%%%%%%%%%%%%%%%%%%

Classical RK4 has tableau

\[
\boxed{
\begin{array}{c|cccc}
0&0&0&0&0\\
1/2&1/2&0&0&0\\
1/2&0&1/2&0&0\\
1&0&0&1&0\\
\hline
&1/6&1/3&1/3&1/6
\end{array}
}
\]

This compact notation describes all stages and final weights.

%%%%%%%%%%%%%%%%%%%%%%%%%%%%%%%%%%%%%%%%%%%%%%%%%%%%%%%%%%%%
\subsection{One-Step Methods}
\label{subsec:one-step-methods}
%%%%%%%%%%%%%%%%%%%%%%%%%%%%%%%%%%%%%%%%%%%%%%%%%%%%%%%%%%%%

A one-step method computes

\[
y_{n+1}
\]

using only information generated from

\[
y_n
\]

during the current step.

Examples include:

\[
\boxed{
\text{Euler},
}
\]

\[
\boxed{
\text{Runge--Kutta methods},
}
\]

and many implicit Runge--Kutta methods.

%%%%%%%%%%%%%%%%%%%%%%%%%%%%%%%%%%%%%%%%%%%%%%%%%%%%%%%%%%%%
\subsection{Linear Multistep Methods}
\label{subsec:linear-multistep-methods}
%%%%%%%%%%%%%%%%%%%%%%%%%%%%%%%%%%%%%%%%%%%%%%%%%%%%%%%%%%%%

A linear multistep method uses several previous values.

A general form is

\[
\boxed{
\sum_{j=0}^{k}
\alpha_j
y_{n+j}
=
h
\sum_{j=0}^{k}
\beta_j
f_{n+j},
}
\]

where

\[
f_{n+j}
=
f(t_{n+j},y_{n+j}).
\]

Previously computed solution history is reused.

%%%%%%%%%%%%%%%%%%%%%%%%%%%%%%%%%%%%%%%%%%%%%%%%%%%%%%%%%%%%
\subsection{Adams--Bashforth Methods}
\label{subsec:adams-bashforth}
%%%%%%%%%%%%%%%%%%%%%%%%%%%%%%%%%%%%%%%%%%%%%%%%%%%%%%%%%%%%

Adams--Bashforth methods are explicit multistep methods.

The two-step Adams--Bashforth method is

\[
\boxed{
y_{n+1}
=
y_n
+
\frac{h}{2}
\left(
3f_n-f_{n-1}
\right).
}
\]

It uses slope information from the two most recent time points.

%%%%%%%%%%%%%%%%%%%%%%%%%%%%%%%%%%%%%%%%%%%%%%%%%%%%%%%%%%%%
\subsection{Three-Step Adams--Bashforth}
\label{subsec:ab3}
%%%%%%%%%%%%%%%%%%%%%%%%%%%%%%%%%%%%%%%%%%%%%%%%%%%%%%%%%%%%

The three-step method is

\[
\boxed{
y_{n+1}
=
y_n
+
\frac{h}{12}
\left(
23f_n
-
16f_{n-1}
+
5f_{n-2}
\right).
}
\]

It is third-order accurate under standard smoothness conditions.

%%%%%%%%%%%%%%%%%%%%%%%%%%%%%%%%%%%%%%%%%%%%%%%%%%%%%%%%%%%%
\subsection{Adams--Moulton Methods}
\label{subsec:adams-moulton}
%%%%%%%%%%%%%%%%%%%%%%%%%%%%%%%%%%%%%%%%%%%%%%%%%%%%%%%%%%%%

Adams--Moulton methods are implicit multistep methods.

The two-point trapezoidal formula can be written

\[
\boxed{
y_{n+1}
=
y_n
+
\frac{h}{2}
(
f_n+f_{n+1}
).
}
\]

A higher-order example is

\[
\boxed{
y_{n+1}
=
y_n
+
\frac{h}{12}
\left(
5f_{n+1}
+
8f_n
-
f_{n-1}
\right).
}
\]

%%%%%%%%%%%%%%%%%%%%%%%%%%%%%%%%%%%%%%%%%%%%%%%%%%%%%%%%%%%%
\subsection{Predictor--Corrector Methods}
\label{subsec:predictor-corrector}
%%%%%%%%%%%%%%%%%%%%%%%%%%%%%%%%%%%%%%%%%%%%%%%%%%%%%%%%%%%%

A predictor--corrector scheme may use an explicit method to estimate

\[
\widetilde{y}_{n+1},
\]

then use that value inside an implicit or more accurate correction.

For example:

\[
\boxed{
\text{Adams--Bashforth predictor}
}
\]

followed by

\[
\boxed{
\text{Adams--Moulton corrector}.
}
\]

This reduces the cost of solving fully implicit equations.

%%%%%%%%%%%%%%%%%%%%%%%%%%%%%%%%%%%%%%%%%%%%%%%%%%%%%%%%%%%%
\subsection{Starting Multistep Methods}
\label{subsec:starting-multistep-methods}
%%%%%%%%%%%%%%%%%%%%%%%%%%%%%%%%%%%%%%%%%%%%%%%%%%%%%%%%%%%%

A \(k\)-step method needs several previous solution values.

At the initial time, these values are not yet available.

Therefore a one-step method such as RK4 may be used to generate startup
values.

%%%%%%%%%%%%%%%%%%%%%%%%%%%%%%%%%%%%%%%%%%%%%%%%%%%%%%%%%%%%
\subsection{Backward Differentiation Formulas}
\label{subsec:backward-differentiation-formulas}
%%%%%%%%%%%%%%%%%%%%%%%%%%%%%%%%%%%%%%%%%%%%%%%%%%%%%%%%%%%%

Backward Differentiation Formulas, abbreviated BDF, approximate the
derivative using backward differences and evaluate the differential
equation at the new state.

Backward Euler is the one-step BDF method.

The two-step BDF formula is

\[
\boxed{
\frac{
3y_{n+1}-4y_n+y_{n-1}
}{
2h
}
=
f(t_{n+1},y_{n+1}).
}
\]

BDF methods are important for stiff differential equations.

%%%%%%%%%%%%%%%%%%%%%%%%%%%%%%%%%%%%%%%%%%%%%%%%%%%%%%%%%%%%
\subsection{Zero-Stability}
\label{subsec:zero-stability}
%%%%%%%%%%%%%%%%%%%%%%%%%%%%%%%%%%%%%%%%%%%%%%%%%%%%%%%%%%%%

Linear multistep methods must satisfy a stability requirement even as

\[
h\to0.
\]

Suppose the first characteristic polynomial is

\[
\boxed{
\rho(\xi)
=
\sum_{j=0}^{k}
\alpha_j\xi^j.
}
\]

The root condition requires:

\[
\boxed{
|\xi_j|\leq1
}
\]

for all roots and that roots with

\[
|\xi_j|=1
\]

be simple.

This is the standard zero-stability condition.

%%%%%%%%%%%%%%%%%%%%%%%%%%%%%%%%%%%%%%%%%%%%%%%%%%%%%%%%%%%%
\subsection{Consistency Plus Zero-Stability}
\label{subsec:dahlquist-equivalence-preview}
%%%%%%%%%%%%%%%%%%%%%%%%%%%%%%%%%%%%%%%%%%%%%%%%%%%%%%%%%%%%

For linear multistep methods, a fundamental result states that

\[
\boxed{
\text{consistency}
+
\text{zero-stability}
\Longleftrightarrow
\text{convergence}.
}
\]

This is the Dahlquist equivalence principle.

It shows why merely matching the differential equation locally is not
enough.

%%%%%%%%%%%%%%%%%%%%%%%%%%%%%%%%%%%%%%%%%%%%%%%%%%%%%%%%%%%%
\subsection{Linear Stability Test Equation}
\label{subsec:ode-test-equation}
%%%%%%%%%%%%%%%%%%%%%%%%%%%%%%%%%%%%%%%%%%%%%%%%%%%%%%%%%%%%

To study numerical stability, consider

\[
\boxed{
y'
=
\lambda y,
}
\]

where

\[
\lambda\in\mathbb{C}.
\]

The exact solution is

\[
\boxed{
y(t)
=
e^{\lambda t}y_0.
}
\]

If

\[
\operatorname{Re}(\lambda)<0,
\]

the exact solution decays.

A good numerical method should reproduce this decay for suitable step
sizes.

%%%%%%%%%%%%%%%%%%%%%%%%%%%%%%%%%%%%%%%%%%%%%%%%%%%%%%%%%%%%
\subsection{Forward Euler Stability Function}
\label{subsec:forward-euler-stability}
%%%%%%%%%%%%%%%%%%%%%%%%%%%%%%%%%%%%%%%%%%%%%%%%%%%%%%%%%%%%

Apply forward Euler:

\[
y_{n+1}
=
y_n+h\lambda y_n.
\]

Therefore,

\[
\boxed{
y_{n+1}
=
(1+h\lambda)y_n.
}
\]

Define

\[
z=h\lambda.
\]

The amplification or stability function is

\[
\boxed{
R(z)=1+z.
}
\]

%%%%%%%%%%%%%%%%%%%%%%%%%%%%%%%%%%%%%%%%%%%%%%%%%%%%%%%%%%%%
\subsection{Absolute Stability Region}
\label{subsec:absolute-stability-region}
%%%%%%%%%%%%%%%%%%%%%%%%%%%%%%%%%%%%%%%%%%%%%%%%%%%%%%%%%%%%

The absolute stability region is

\[
\boxed{
\mathcal{S}
=
\{
z\in\mathbb{C}
:
|R(z)|\leq1
\}.
}
\]

For forward Euler,

\[
\boxed{
|1+z|
\leq1.
}
\]

This is the disk centered at

\[
-1
\]

with radius

\[
1.
\]

%%%%%%%%%%%%%%%%%%%%%%%%%%%%%%%%%%%%%%%%%%%%%%%%%%%%%%%%%%%%
\subsection{Worked Example: Euler Stability Restriction}
\label{subsec:worked-euler-stability}
%%%%%%%%%%%%%%%%%%%%%%%%%%%%%%%%%%%%%%%%%%%%%%%%%%%%%%%%%%%%

\begin{workedexamplebox}{Stable Step for Rapid Decay}

Consider

\[
y'=-100y.
\]

Forward Euler gives

\[
y_{n+1}
=
(1-100h)y_n.
\]

For absolute stability,

\[
|1-100h|
\leq1.
\]

Therefore,

\[
-1
\leq
1-100h
\leq
1.
\]

This yields

\[
0
\leq
h
\leq
0.02.
\]

\begin{answerbox}
\[
\boxed{
h\leq0.02
}
\]

is required for forward Euler absolute stability.
\end{answerbox}

\end{workedexamplebox}

%%%%%%%%%%%%%%%%%%%%%%%%%%%%%%%%%%%%%%%%%%%%%%%%%%%%%%%%%%%%
\subsection{Backward Euler Stability Function}
\label{subsec:backward-euler-stability}
%%%%%%%%%%%%%%%%%%%%%%%%%%%%%%%%%%%%%%%%%%%%%%%%%%%%%%%%%%%%

Apply backward Euler to

\[
y'=\lambda y:
\]

\[
y_{n+1}
=
y_n+h\lambda y_{n+1}.
\]

Hence,

\[
(1-h\lambda)y_{n+1}
=
y_n.
\]

Therefore,

\[
\boxed{
R(z)
=
\frac1{1-z}.
}
\]

%%%%%%%%%%%%%%%%%%%%%%%%%%%%%%%%%%%%%%%%%%%%%%%%%%%%%%%%%%%%
\subsection{A-Stability}
\label{subsec:a-stability}
%%%%%%%%%%%%%%%%%%%%%%%%%%%%%%%%%%%%%%%%%%%%%%%%%%%%%%%%%%%%

A method is A-stable if its absolute stability region contains the entire
left half-plane:

\[
\boxed{
\operatorname{Re}(z)\leq0.
}
\]

Backward Euler is A-stable.

Thus all exponentially decaying linear test equations remain stable for
every positive step size.

Accuracy may still require small steps.

%%%%%%%%%%%%%%%%%%%%%%%%%%%%%%%%%%%%%%%%%%%%%%%%%%%%%%%%%%%%
\subsection{Trapezoidal Stability}
\label{subsec:trapezoidal-stability}
%%%%%%%%%%%%%%%%%%%%%%%%%%%%%%%%%%%%%%%%%%%%%%%%%%%%%%%%%%%%

For the test equation,

\[
y_{n+1}
=
y_n
+
\frac{z}{2}
(
y_n+y_{n+1}
).
\]

Therefore,

\[
\boxed{
R(z)
=
\frac{
1+z/2
}{
1-z/2
}.
}
\]

The trapezoidal method is A-stable.

However, for very large negative \(z\),

\[
R(z)\to-1,
\]

so strongly damped modes may produce persistent numerical oscillation.

%%%%%%%%%%%%%%%%%%%%%%%%%%%%%%%%%%%%%%%%%%%%%%%%%%%%%%%%%%%%
\subsection{L-Stability}
\label{subsec:l-stability}
%%%%%%%%%%%%%%%%%%%%%%%%%%%%%%%%%%%%%%%%%%%%%%%%%%%%%%%%%%%%

An A-stable method is L-stable if additionally

\[
\boxed{
R(z)\to0
}
\]

as

\[
z\to-\infty
\]

along the negative real axis or, more generally, for strongly stable
modes in the relevant limit.

Backward Euler is L-stable.

The trapezoidal method is not.

L-stability is useful when rapidly decaying modes should be strongly
damped numerically.

%%%%%%%%%%%%%%%%%%%%%%%%%%%%%%%%%%%%%%%%%%%%%%%%%%%%%%%%%%%%
\subsection{Stiff Differential Equations}
\label{subsec:stiff-odes}
%%%%%%%%%%%%%%%%%%%%%%%%%%%%%%%%%%%%%%%%%%%%%%%%%%%%%%%%%%%%

A differential equation is called stiff when stable numerical integration
requires an extremely small time step for some methods even though the
solution itself changes slowly over much of the interval.

There is no single universal definition of stiffness.

A common signature is the coexistence of widely separated time scales.

%%%%%%%%%%%%%%%%%%%%%%%%%%%%%%%%%%%%%%%%%%%%%%%%%%%%%%%%%%%%
\subsection{Simple Stiff Example}
\label{subsec:simple-stiff-example}
%%%%%%%%%%%%%%%%%%%%%%%%%%%%%%%%%%%%%%%%%%%%%%%%%%%%%%%%%%%%

Consider

\[
\boxed{
y'
=
-1000
(
y-\cos t
)
-
\sin t.
}
\]

The exact solution contains a rapidly decaying transient together with a
slow component related to

\[
\cos t.
\]

An explicit method may be forced to resolve the fast stable mode even after
it has become negligible.

%%%%%%%%%%%%%%%%%%%%%%%%%%%%%%%%%%%%%%%%%%%%%%%%%%%%%%%%%%%%
\subsection{Why Explicit Euler Struggles With Stiffness}
\label{subsec:euler-stiffness}
%%%%%%%%%%%%%%%%%%%%%%%%%%%%%%%%%%%%%%%%%%%%%%%%%%%%%%%%%%%%

For

\[
y'=\lambda y
\]

with large negative \(\lambda\), forward Euler requires

\[
\boxed{
|1+h\lambda|\leq1.
}
\]

Thus

\[
h
\]

must be very small.

The restriction arises from numerical stability rather than from the
smoothness of the slowly varying physical solution.

%%%%%%%%%%%%%%%%%%%%%%%%%%%%%%%%%%%%%%%%%%%%%%%%%%%%%%%%%%%%
\subsection{Implicit Methods for Stiff Problems}
\label{subsec:implicit-stiff-methods}
%%%%%%%%%%%%%%%%%%%%%%%%%%%%%%%%%%%%%%%%%%%%%%%%%%%%%%%%%%%%

A-stable or L-stable implicit methods permit much larger stable steps for
rapidly decaying modes.

Common stiff methods include:

\[
\boxed{
\text{backward Euler},
}
\]

\[
\boxed{
\text{BDF methods},
}
\]

and various

\[
\boxed{
\text{implicit Runge--Kutta methods}.
}
\]

The tradeoff is that nonlinear or linear algebraic systems must be solved
at each step.

%%%%%%%%%%%%%%%%%%%%%%%%%%%%%%%%%%%%%%%%%%%%%%%%%%%%%%%%%%%%
\subsection{Systems of ODEs}
\label{subsec:numerical-ode-systems}
%%%%%%%%%%%%%%%%%%%%%%%%%%%%%%%%%%%%%%%%%%%%%%%%%%%%%%%%%%%%

A system has form

\[
\boxed{
\mathbf{y}'(t)
=
\mathbf{f}
(
t,\mathbf{y}(t)
),
}
\]

where

\[
\mathbf{y}(t)
\in
\mathbb{R}^m.
\]

Euler becomes

\[
\boxed{
\mathbf{y}_{n+1}
=
\mathbf{y}_n
+
h
\mathbf{f}
(
t_n,\mathbf{y}_n
).
}
\]

Runge--Kutta methods extend componentwise in vector form.

%%%%%%%%%%%%%%%%%%%%%%%%%%%%%%%%%%%%%%%%%%%%%%%%%%%%%%%%%%%%
\subsection{Higher-Order ODEs as First-Order Systems}
\label{subsec:higher-order-to-first-order}
%%%%%%%%%%%%%%%%%%%%%%%%%%%%%%%%%%%%%%%%%%%%%%%%%%%%%%%%%%%%

Consider

\[
\boxed{
y''=F(t,y,y').
}
\]

Define

\[
\boxed{
y_1=y,
}
\]

and

\[
\boxed{
y_2=y'.
}
\]

Then

\[
\boxed{
y_1'=y_2,
}
\]

and

\[
\boxed{
y_2'
=
F(t,y_1,y_2).
}
\]

Thus numerical first-order methods can solve higher-order ODEs after state
augmentation.

%%%%%%%%%%%%%%%%%%%%%%%%%%%%%%%%%%%%%%%%%%%%%%%%%%%%%%%%%%%%
\subsection{Worked Example: Harmonic Oscillator}
\label{subsec:worked-harmonic-oscillator-system}
%%%%%%%%%%%%%%%%%%%%%%%%%%%%%%%%%%%%%%%%%%%%%%%%%%%%%%%%%%%%

\begin{workedexamplebox}{Converting a Second-Order ODE}

Consider

\[
y''+y=0.
\]

Define

\[
y_1=y,
\qquad
y_2=y'.
\]

Then

\[
y_1'=y_2,
\]

and

\[
y_2'=-y_1.
\]

Thus,

\begin{answerbox}
\[
\boxed{
\begin{pmatrix}
y_1'\\
y_2'
\end{pmatrix}
=
\begin{pmatrix}
y_2\\
-y_1
\end{pmatrix}.
}
\]
\end{answerbox}

This system can be integrated using any standard numerical ODE solver.

\end{workedexamplebox}

%%%%%%%%%%%%%%%%%%%%%%%%%%%%%%%%%%%%%%%%%%%%%%%%%%%%%%%%%%%%
\subsection{Jacobian Matrix}
\label{subsec:ode-jacobian}
%%%%%%%%%%%%%%%%%%%%%%%%%%%%%%%%%%%%%%%%%%%%%%%%%%%%%%%%%%%%

For a system

\[
\mathbf{f}
(
t,\mathbf{y}
),
\]

the state Jacobian is

\[
\boxed{
J
=
\frac{
\partial\mathbf{f}
}{
\partial\mathbf{y}
}.
}
\]

Its entries are

\[
\boxed{
J_{ij}
=
\frac{
\partial f_i
}{
\partial y_j
}.
}
\]

Implicit methods often require this Jacobian or an approximation to it.

%%%%%%%%%%%%%%%%%%%%%%%%%%%%%%%%%%%%%%%%%%%%%%%%%%%%%%%%%%%%
\subsection{Newton Solve for Backward Euler Systems}
\label{subsec:newton-backward-euler-system}
%%%%%%%%%%%%%%%%%%%%%%%%%%%%%%%%%%%%%%%%%%%%%%%%%%%%%%%%%%%%

Backward Euler requires solving

\[
\boxed{
\mathbf{G}(\mathbf{z})
=
\mathbf{z}
-
\mathbf{y}_n
-
h
\mathbf{f}
(
t_{n+1},\mathbf{z}
)
=
0.
}
\]

Its Jacobian is

\[
\boxed{
G'(\mathbf{z})
=
I
-
h
J_f
(
t_{n+1},\mathbf{z}
).
}
\]

Newton's method solves

\[
\boxed{
\left[
I-hJ_f
\right]
\Delta
=
-\mathbf{G},
}
\]

then updates

\[
\boxed{
\mathbf{z}
\leftarrow
\mathbf{z}+\Delta.
}
\]

%%%%%%%%%%%%%%%%%%%%%%%%%%%%%%%%%%%%%%%%%%%%%%%%%%%%%%%%%%%%
\subsection{Cost of Implicit Integration}
\label{subsec:implicit-integration-cost}
%%%%%%%%%%%%%%%%%%%%%%%%%%%%%%%%%%%%%%%%%%%%%%%%%%%%%%%%%%%%

An implicit step may require:

\[
\boxed{
\text{function evaluations},
}
\]

\[
\boxed{
\text{Jacobian evaluation},
}
\]

\[
\boxed{
\text{matrix factorization},
}
\]

and

\[
\boxed{
\text{multiple Newton iterations}.
}
\]

Thus one implicit step is usually more expensive than one explicit step.

For stiff systems, however, far fewer time steps may be needed.

%%%%%%%%%%%%%%%%%%%%%%%%%%%%%%%%%%%%%%%%%%%%%%%%%%%%%%%%%%%%
\subsection{Jacobian Reuse}
\label{subsec:ode-jacobian-reuse}
%%%%%%%%%%%%%%%%%%%%%%%%%%%%%%%%%%%%%%%%%%%%%%%%%%%%%%%%%%%%

If the Jacobian changes slowly, an implicit solver may reuse a previous
factorization for several nonlinear iterations or time steps.

This can substantially reduce computational cost.

The solver must detect when the old Jacobian approximation has become
inadequate.

%%%%%%%%%%%%%%%%%%%%%%%%%%%%%%%%%%%%%%%%%%%%%%%%%%%%%%%%%%%%
\subsection{Adaptive Time Stepping}
\label{subsec:adaptive-time-stepping}
%%%%%%%%%%%%%%%%%%%%%%%%%%%%%%%%%%%%%%%%%%%%%%%%%%%%%%%%%%%%

A fixed step size treats every part of the trajectory equally.

But some intervals may be smooth while others change rapidly.

Adaptive methods choose

\[
\boxed{
h_n
}
\]

dynamically based on estimated local error.

The goals are:

\[
\boxed{
\text{accuracy}
}
\]

and

\[
\boxed{
\text{computational efficiency}.
}
\]

%%%%%%%%%%%%%%%%%%%%%%%%%%%%%%%%%%%%%%%%%%%%%%%%%%%%%%%%%%%%
\subsection{Embedded Runge--Kutta Methods}
\label{subsec:embedded-runge-kutta}
%%%%%%%%%%%%%%%%%%%%%%%%%%%%%%%%%%%%%%%%%%%%%%%%%%%%%%%%%%%%

An embedded Runge--Kutta pair uses the same stage evaluations to construct
two approximations of different orders.

For example,

\[
\boxed{
y_{n+1}^{[p]}
}
\]

and

\[
\boxed{
y_{n+1}^{[p+1]}.
}
\]

Their difference estimates local error:

\[
\boxed{
e_{n+1}
\approx
y_{n+1}^{[p+1]}
-
y_{n+1}^{[p]}.
}
\]

%%%%%%%%%%%%%%%%%%%%%%%%%%%%%%%%%%%%%%%%%%%%%%%%%%%%%%%%%%%%
\subsection{Scaled Error}
\label{subsec:ode-scaled-error}
%%%%%%%%%%%%%%%%%%%%%%%%%%%%%%%%%%%%%%%%%%%%%%%%%%%%%%%%%%%%

Absolute error alone is not appropriate for state components with
different scales.

A componentwise tolerance may use

\[
\boxed{
\mathrm{scale}_i
=
\mathrm{atol}_i
+
\mathrm{rtol}
\max
\{
|y_{n,i}|,
|y_{n+1,i}|
\}.
}
\]

Then form normalized errors such as

\[
\boxed{
\frac{
e_i
}{
\mathrm{scale}_i
}.
}
\]

%%%%%%%%%%%%%%%%%%%%%%%%%%%%%%%%%%%%%%%%%%%%%%%%%%%%%%%%%%%%
\subsection{Absolute and Relative Tolerance}
\label{subsec:ode-atol-rtol}
%%%%%%%%%%%%%%%%%%%%%%%%%%%%%%%%%%%%%%%%%%%%%%%%%%%%%%%%%%%%

The absolute tolerance

\[
\boxed{
\mathrm{atol}
}
\]

controls error for quantities near zero.

The relative tolerance

\[
\boxed{
\mathrm{rtol}
}
\]

controls error relative to the magnitude of the solution.

Using both is generally more meaningful than using either alone.

%%%%%%%%%%%%%%%%%%%%%%%%%%%%%%%%%%%%%%%%%%%%%%%%%%%%%%%%%%%%
\subsection{Step Acceptance}
\label{subsec:ode-step-acceptance}
%%%%%%%%%%%%%%%%%%%%%%%%%%%%%%%%%%%%%%%%%%%%%%%%%%%%%%%%%%%%

Let

\[
E_n
\]

be a normalized local-error estimate.

A typical decision is:

\[
\boxed{
E_n\leq1
\Rightarrow
\text{accept the step},
}
\]

and

\[
\boxed{
E_n>1
\Rightarrow
\text{reject and retry with smaller }h.
}
\]

%%%%%%%%%%%%%%%%%%%%%%%%%%%%%%%%%%%%%%%%%%%%%%%%%%%%%%%%%%%%
\subsection{Adaptive Step-Size Formula}
\label{subsec:adaptive-step-formula}
%%%%%%%%%%%%%%%%%%%%%%%%%%%%%%%%%%%%%%%%%%%%%%%%%%%%%%%%%%%%

Suppose local error behaves approximately like

\[
C h^{p+1}.
\]

Then a new step may be selected using

\[
\boxed{
h_{\mathrm{new}}
=
s
h
E_n^{-1/(p+1)},
}
\]

where

\[
0<s<1
\]

is a safety factor.

Practical solvers also limit how rapidly the step size may increase or
decrease.

%%%%%%%%%%%%%%%%%%%%%%%%%%%%%%%%%%%%%%%%%%%%%%%%%%%%%%%%%%%%
\subsection{Rejected Steps}
\label{subsec:ode-rejected-steps}
%%%%%%%%%%%%%%%%%%%%%%%%%%%%%%%%%%%%%%%%%%%%%%%%%%%%%%%%%%%%

If the estimated error is too large, the step is rejected.

The numerical state returns to the beginning of the failed step.

The solver then retries using a smaller step size.

Rejected steps are part of normal adaptive integration.

%%%%%%%%%%%%%%%%%%%%%%%%%%%%%%%%%%%%%%%%%%%%%%%%%%%%%%%%%%%%
\subsection{Dense Output}
\label{subsec:ode-dense-output}
%%%%%%%%%%%%%%%%%%%%%%%%%%%%%%%%%%%%%%%%%%%%%%%%%%%%%%%%%%%%

A numerical solver computes values at discrete internal time points.

Applications may require the solution at arbitrary intermediate times.

Dense output constructs an interpolation formula that is consistent with
the time-stepping method.

This is more accurate than arbitrary interpolation between stored points.

%%%%%%%%%%%%%%%%%%%%%%%%%%%%%%%%%%%%%%%%%%%%%%%%%%%%%%%%%%%%
\subsection{Event Detection}
\label{subsec:ode-event-detection}
%%%%%%%%%%%%%%%%%%%%%%%%%%%%%%%%%%%%%%%%%%%%%%%%%%%%%%%%%%%%

An event occurs when a function

\[
\boxed{
g(t,y(t))
}
\]

crosses zero.

Examples include:

\[
\boxed{
\text{impact with the ground},
}
\]

\[
\boxed{
\text{threshold crossing},
}
\]

or

\[
\boxed{
\text{population extinction}.
}
\]

The solver brackets and locates the event time using interpolation or
root-finding.

%%%%%%%%%%%%%%%%%%%%%%%%%%%%%%%%%%%%%%%%%%%%%%%%%%%%%%%%%%%%
\subsection{Terminal and Nonterminal Events}
\label{subsec:terminal-nonterminal-events}
%%%%%%%%%%%%%%%%%%%%%%%%%%%%%%%%%%%%%%%%%%%%%%%%%%%%%%%%%%%%

A terminal event stops integration.

A nonterminal event is recorded, but integration continues.

Event functions may also specify which crossing direction matters:

\[
\boxed{
\text{negative to positive},
}
\]

\[
\boxed{
\text{positive to negative},
}
\]

or either.

%%%%%%%%%%%%%%%%%%%%%%%%%%%%%%%%%%%%%%%%%%%%%%%%%%%%%%%%%%%%
\subsection{Conservation Laws}
\label{subsec:numerical-ode-conservation}
%%%%%%%%%%%%%%%%%%%%%%%%%%%%%%%%%%%%%%%%%%%%%%%%%%%%%%%%%%%%

A differential equation may conserve a quantity

\[
\boxed{
I(y(t))
=
\text{constant}.
}
\]

A general-purpose numerical method may not preserve this invariant
exactly.

The numerical invariant may slowly drift.

This can become important in long simulations.

%%%%%%%%%%%%%%%%%%%%%%%%%%%%%%%%%%%%%%%%%%%%%%%%%%%%%%%%%%%%
\subsection{Energy Drift}
\label{subsec:energy-drift}
%%%%%%%%%%%%%%%%%%%%%%%%%%%%%%%%%%%%%%%%%%%%%%%%%%%%%%%%%%%%

For mechanical systems, the true energy may satisfy

\[
\boxed{
H(q(t),p(t))
=
H(q(0),p(0)).
}
\]

Standard numerical methods may introduce artificial energy growth or
decay.

Long-time qualitative behavior can therefore depend strongly on method
choice.

%%%%%%%%%%%%%%%%%%%%%%%%%%%%%%%%%%%%%%%%%%%%%%%%%%%%%%%%%%%%
\subsection{Symplectic Integrators}
\label{subsec:symplectic-integrators}
%%%%%%%%%%%%%%%%%%%%%%%%%%%%%%%%%%%%%%%%%%%%%%%%%%%%%%%%%%%%

Symplectic integrators are designed for Hamiltonian systems.

They preserve the underlying symplectic geometry rather than generally
preserving energy exactly.

Over long times, they often reproduce qualitative mechanical behavior much
better than generic methods of comparable order.

%%%%%%%%%%%%%%%%%%%%%%%%%%%%%%%%%%%%%%%%%%%%%%%%%%%%%%%%%%%%
\subsection{Symplectic Euler Preview}
\label{subsec:symplectic-euler}
%%%%%%%%%%%%%%%%%%%%%%%%%%%%%%%%%%%%%%%%%%%%%%%%%%%%%%%%%%%%

For a separable Hamiltonian system

\[
\boxed{
\dot{q}
=
\frac{\partial H}{\partial p},
\qquad
\dot{p}
=
-
\frac{\partial H}{\partial q},
}
\]

a symplectic Euler variant updates one variable implicitly and the other
explicitly.

For example,

\[
\boxed{
p_{n+1}
=
p_n
-
h
\nabla V(q_n),
}
\]

followed by

\[
\boxed{
q_{n+1}
=
q_n
+
h
M^{-1}p_{n+1}.
}
\]

%%%%%%%%%%%%%%%%%%%%%%%%%%%%%%%%%%%%%%%%%%%%%%%%%%%%%%%%%%%%
\subsection{Positivity Preservation}
\label{subsec:ode-positivity-preservation}
%%%%%%%%%%%%%%%%%%%%%%%%%%%%%%%%%%%%%%%%%%%%%%%%%%%%%%%%%%%%

Some physical variables must satisfy

\[
\boxed{
y(t)\geq0.
}
\]

Examples include:

\[
\boxed{
\text{population},
}
\]

\[
\boxed{
\text{concentration},
}
\]

and

\[
\boxed{
\text{probability}.
}
\]

A numerical method may produce small negative values even when the exact
solution remains nonnegative.

Method and step-size choices should respect important physical invariants
when possible.

%%%%%%%%%%%%%%%%%%%%%%%%%%%%%%%%%%%%%%%%%%%%%%%%%%%%%%%%%%%%
\subsection{Population Conservation}
\label{subsec:population-conservation}
%%%%%%%%%%%%%%%%%%%%%%%%%%%%%%%%%%%%%%%%%%%%%%%%%%%%%%%%%%%%

Compartmental models may satisfy

\[
\boxed{
S+E+I+R+D=N.
}
\]

A numerical scheme may preserve this exactly, approximately, or poorly.

Checking conserved totals is a useful verification tool.

%%%%%%%%%%%%%%%%%%%%%%%%%%%%%%%%%%%%%%%%%%%%%%%%%%%%%%%%%%%%
\subsection{Boundary-Value Problems}
\label{subsec:numerical-ode-boundary-value}
%%%%%%%%%%%%%%%%%%%%%%%%%%%%%%%%%%%%%%%%%%%%%%%%%%%%%%%%%%%%

An ODE boundary-value problem specifies conditions at more than one point.

For example,

\[
\boxed{
y''=f(t,y,y'),
}
\]

with

\[
\boxed{
y(a)=\alpha,
\qquad
y(b)=\beta.
}
\]

Unlike an initial-value problem, the required initial derivative may be
unknown.

%%%%%%%%%%%%%%%%%%%%%%%%%%%%%%%%%%%%%%%%%%%%%%%%%%%%%%%%%%%%
\subsection{Shooting Method}
\label{subsec:ode-shooting-method}
%%%%%%%%%%%%%%%%%%%%%%%%%%%%%%%%%%%%%%%%%%%%%%%%%%%%%%%%%%%%

For

\[
y''=f(t,y,y'),
\]

with

\[
y(a)=\alpha,
\qquad
y(b)=\beta,
\]

guess

\[
\boxed{
s=y'(a).
}
\]

Solve the initial-value problem using this guess.

Then define terminal mismatch

\[
\boxed{
F(s)
=
y(b;s)-\beta.
}
\]

Adjust \(s\) using root-finding until

\[
\boxed{
F(s)=0.
}
\]

%%%%%%%%%%%%%%%%%%%%%%%%%%%%%%%%%%%%%%%%%%%%%%%%%%%%%%%%%%%%
\subsection{Worked Example: Shooting Formulation}
\label{subsec:worked-shooting-formulation}
%%%%%%%%%%%%%%%%%%%%%%%%%%%%%%%%%%%%%%%%%%%%%%%%%%%%%%%%%%%%

\begin{workedexamplebox}{Turning a BVP Into Root Finding}

Suppose

\[
y''=-y,
\]

with

\[
y(0)=0,
\qquad
y(1)=1.
\]

Let

\[
s=y'(0)
\]

be unknown.

Solve the IVP

\[
y(0)=0,
\qquad
y'(0)=s.
\]

Define

\[
F(s)
=
y(1;s)-1.
\]

The boundary-value problem becomes

\[
\boxed{
F(s)=0.
}
\]

Thus a numerical ODE solver is nested inside a numerical root solver.

\end{workedexamplebox}

%%%%%%%%%%%%%%%%%%%%%%%%%%%%%%%%%%%%%%%%%%%%%%%%%%%%%%%%%%%%
\subsection{Multiple Shooting}
\label{subsec:ode-multiple-shooting}
%%%%%%%%%%%%%%%%%%%%%%%%%%%%%%%%%%%%%%%%%%%%%%%%%%%%%%%%%%%%

Single shooting can become unstable when the IVP is highly sensitive to
the unknown initial conditions.

Multiple shooting divides the interval into subintervals.

Separate initial states are introduced on each segment and continuity
constraints connect them.

This can improve numerical conditioning.

%%%%%%%%%%%%%%%%%%%%%%%%%%%%%%%%%%%%%%%%%%%%%%%%%%%%%%%%%%%%
\subsection{Finite Differences for Boundary-Value Problems}
\label{subsec:ode-bvp-finite-differences}
%%%%%%%%%%%%%%%%%%%%%%%%%%%%%%%%%%%%%%%%%%%%%%%%%%%%%%%%%%%%

For

\[
y''(x)=f(x,y),
\]

use

\[
\boxed{
y''(x_i)
\approx
\frac{
y_{i+1}-2y_i+y_{i-1}
}{
h^2
}.
}
\]

Then the differential equation becomes a system of algebraic equations for

\[
y_1,\ldots,y_{N-1}.
\]

Boundary values are inserted directly.

%%%%%%%%%%%%%%%%%%%%%%%%%%%%%%%%%%%%%%%%%%%%%%%%%%%%%%%%%%%%
\subsection{Nonlinear Boundary-Value Systems}
\label{subsec:nonlinear-bvp-systems}
%%%%%%%%%%%%%%%%%%%%%%%%%%%%%%%%%%%%%%%%%%%%%%%%%%%%%%%%%%%%

If the differential equation is nonlinear, finite-difference
discretization produces a nonlinear algebraic system:

\[
\boxed{
F(Y)=0.
}
\]

Newton's method can then be applied to the entire discrete system.

This again demonstrates the connection between differential equations,
numerical linear algebra, and nonlinear equation solving.

%%%%%%%%%%%%%%%%%%%%%%%%%%%%%%%%%%%%%%%%%%%%%%%%%%%%%%%%%%%%
\subsection{Method of Lines Preview}
\label{subsec:method-of-lines-preview}
%%%%%%%%%%%%%%%%%%%%%%%%%%%%%%%%%%%%%%%%%%%%%%%%%%%%%%%%%%%%

A time-dependent PDE can often be discretized only in space.

The spatial derivatives are replaced by finite-dimensional
approximations.

The result is a large system of ODEs:

\[
\boxed{
\mathbf{u}'(t)
=
F(t,\mathbf{u}(t)).
}
\]

An ODE solver then integrates the semi-discrete system in time.

This is the method of lines.

%%%%%%%%%%%%%%%%%%%%%%%%%%%%%%%%%%%%%%%%%%%%%%%%%%%%%%%%%%%%
\subsection{Error Propagation}
\label{subsec:ode-error-propagation}
%%%%%%%%%%%%%%%%%%%%%%%%%%%%%%%%%%%%%%%%%%%%%%%%%%%%%%%%%%%%

Suppose a numerical perturbation at one step is

\[
\delta_n.
\]

The dynamics and numerical method propagate that perturbation into future
steps.

For nonlinear equations, sensitivity may depend on derivatives such as

\[
\boxed{
\frac{\partial f}{\partial y}.
}
\]

Stable methods control the amplification of these perturbations.

%%%%%%%%%%%%%%%%%%%%%%%%%%%%%%%%%%%%%%%%%%%%%%%%%%%%%%%%%%%%
\subsection{Lipschitz Condition}
\label{subsec:ode-lipschitz-numerical}
%%%%%%%%%%%%%%%%%%%%%%%%%%%%%%%%%%%%%%%%%%%%%%%%%%%%%%%%%%%%

Suppose

\[
\boxed{
|f(t,y)-f(t,z)|
\leq
L|y-z|.
}
\]

This Lipschitz condition helps control solution sensitivity and numerical
error growth.

It also appears in existence and uniqueness theory for ODEs.

%%%%%%%%%%%%%%%%%%%%%%%%%%%%%%%%%%%%%%%%%%%%%%%%%%%%%%%%%%%%
\subsection{Discrete Gronwall-Type Error Growth}
\label{subsec:discrete-gronwall-ode}
%%%%%%%%%%%%%%%%%%%%%%%%%%%%%%%%%%%%%%%%%%%%%%%%%%%%%%%%%%%%

A typical error recurrence has form

\[
\boxed{
|e_{n+1}|
\leq
(1+Lh)|e_n|
+
C h^{p+1}.
}
\]

Repeated application gives a bound resembling

\[
\boxed{
|e_n|
\leq
C'
h^p
e^{L(t_n-t_0)}.
}
\]

This explains how local error becomes global error.

%%%%%%%%%%%%%%%%%%%%%%%%%%%%%%%%%%%%%%%%%%%%%%%%%%%%%%%%%%%%
\subsection{Convergence Study}
\label{subsec:ode-convergence-study}
%%%%%%%%%%%%%%%%%%%%%%%%%%%%%%%%%%%%%%%%%%%%%%%%%%%%%%%%%%%%

To test a solver:

\begin{enumerate}
    \item solve with step size \(h\);
    \item solve with \(h/2\);
    \item solve with \(h/4\);
    \item compare against an exact or highly accurate reference solution;
    \item estimate the observed convergence order.
\end{enumerate}

For an order-\(p\) method,

\[
\boxed{
E(h)
\approx
C h^p.
}
\]

%%%%%%%%%%%%%%%%%%%%%%%%%%%%%%%%%%%%%%%%%%%%%%%%%%%%%%%%%%%%
\subsection{Estimating ODE Convergence Order}
\label{subsec:estimating-ode-order}
%%%%%%%%%%%%%%%%%%%%%%%%%%%%%%%%%%%%%%%%%%%%%%%%%%%%%%%%%%%%

If

\[
E(h)
\]

and

\[
E(h/2)
\]

are known, then

\[
\boxed{
p
\approx
\log_2
\left(
\frac{
E(h)
}{
E(h/2)
}
\right).
}
\]

For example, if error decreases by a factor of approximately \(16\), then

\[
\boxed{
p\approx4.
}
\]

This is consistent with fourth-order convergence.

%%%%%%%%%%%%%%%%%%%%%%%%%%%%%%%%%%%%%%%%%%%%%%%%%%%%%%%%%%%%
\subsection{Step Size and Computational Cost}
\label{subsec:ode-step-size-cost}
%%%%%%%%%%%%%%%%%%%%%%%%%%%%%%%%%%%%%%%%%%%%%%%%%%%%%%%%%%%%

For a fixed interval of length \(T\), using step size \(h\) requires
approximately

\[
\boxed{
N
=
\frac{T}{h}
}
\]

steps.

Halving \(h\) approximately doubles the number of steps.

Higher-order methods may use more function evaluations per step but require
fewer steps for a given accuracy.

%%%%%%%%%%%%%%%%%%%%%%%%%%%%%%%%%%%%%%%%%%%%%%%%%%%%%%%%%%%%
\subsection{Function Evaluation Cost}
\label{subsec:ode-function-evaluation-cost}
%%%%%%%%%%%%%%%%%%%%%%%%%%%%%%%%%%%%%%%%%%%%%%%%%%%%%%%%%%%%

For many simulation problems, evaluating

\[
f(t,y)
\]

is expensive.

Forward Euler requires one new evaluation per step.

Classical RK4 requires four stage evaluations per step.

An adaptive higher-order method may still be cheaper overall because it
can take much larger time steps.

%%%%%%%%%%%%%%%%%%%%%%%%%%%%%%%%%%%%%%%%%%%%%%%%%%%%%%%%%%%%
\subsection{Stiffness Detection}
\label{subsec:ode-stiffness-detection}
%%%%%%%%%%%%%%%%%%%%%%%%%%%%%%%%%%%%%%%%%%%%%%%%%%%%%%%%%%%%

Practical solvers may detect stiffness when:

\[
\boxed{
\text{explicit steps repeatedly fail},
}
\]

\[
\boxed{
\text{step sizes become extremely small},
}
\]

or when estimated Jacobian eigenvalues indicate rapid stable modes.

Some software can switch automatically between nonstiff and stiff
integration strategies.

%%%%%%%%%%%%%%%%%%%%%%%%%%%%%%%%%%%%%%%%%%%%%%%%%%%%%%%%%%%%
\subsection{Nonstiff Method Selection}
\label{subsec:nonstiff-method-selection}
%%%%%%%%%%%%%%%%%%%%%%%%%%%%%%%%%%%%%%%%%%%%%%%%%%%%%%%%%%%%

For smooth nonstiff initial-value problems, adaptive explicit Runge--Kutta
methods are often effective.

Advantages include:

\[
\boxed{
\text{no nonlinear solve},
}
\]

\[
\boxed{
\text{good high-order accuracy},
}
\]

and

\[
\boxed{
\text{straightforward adaptive stepping}.
}
\]

%%%%%%%%%%%%%%%%%%%%%%%%%%%%%%%%%%%%%%%%%%%%%%%%%%%%%%%%%%%%
\subsection{Stiff Method Selection}
\label{subsec:stiff-method-selection}
%%%%%%%%%%%%%%%%%%%%%%%%%%%%%%%%%%%%%%%%%%%%%%%%%%%%%%%%%%%%

For stiff problems, consider methods with strong stability properties such
as:

\[
\boxed{
\text{BDF},
}
\]

\[
\boxed{
\text{implicit Runge--Kutta},
}
\]

or other specialized stiff integrators.

Providing an accurate Jacobian can substantially improve solver
performance.

%%%%%%%%%%%%%%%%%%%%%%%%%%%%%%%%%%%%%%%%%%%%%%%%%%%%%%%%%%%%
\subsection{Numerical ODE Workflow}
\label{subsec:numerical-ode-workflow}
%%%%%%%%%%%%%%%%%%%%%%%%%%%%%%%%%%%%%%%%%%%%%%%%%%%%%%%%%%%%

A practical workflow is:

\begin{enumerate}
    \item write the ODE or ODE system clearly;
    \item specify initial conditions;
    \item convert higher-order equations to first-order form;
    \item determine the integration interval;
    \item identify important physical constraints and invariants;
    \item assess whether stiffness is likely;
    \item choose an explicit or implicit solver;
    \item choose absolute and relative tolerances;
    \item supply a Jacobian if appropriate;
    \item configure event functions if needed;
    \item compute the numerical trajectory;
    \item inspect rejected steps and step-size behavior;
    \item verify conserved quantities or known qualitative behavior;
    \item repeat with tighter tolerances;
    \item perform a convergence study;
    \item compare with exact solutions or independent methods when
          available;
    \item confirm that the numerical solution is sufficiently accurate for
          the intended application.
\end{enumerate}

%%%%%%%%%%%%%%%%%%%%%%%%%%%%%%%%%%%%%%%%%%%%%%%%%%%%%%%%%%%%
\subsection{Choosing a Numerical ODE Method}
\label{subsec:choosing-numerical-ode-method}
%%%%%%%%%%%%%%%%%%%%%%%%%%%%%%%%%%%%%%%%%%%%%%%%%%%%%%%%%%%%

For introductory hand calculations, use:

\[
\boxed{
\text{forward Euler}
}
\]

or a low-order Runge--Kutta method.

For general smooth nonstiff problems, use:

\[
\boxed{
\text{adaptive explicit Runge--Kutta}.
}
\]

For stiff systems, consider:

\[
\boxed{
\text{implicit or BDF-type methods}.
}
\]

For long-time Hamiltonian dynamics, consider:

\[
\boxed{
\text{structure-preserving or symplectic methods}.
}
\]

For boundary-value problems, consider:

\[
\boxed{
\text{shooting}
}
\]

or

\[
\boxed{
\text{finite-difference/collocation approaches}.
}
\]

%%%%%%%%%%%%%%%%%%%%%%%%%%%%%%%%%%%%%%%%%%%%%%%%%%%%%%%%%%%%
\subsection{Common Errors}
\label{subsec:numerical-ode-common-errors}
%%%%%%%%%%%%%%%%%%%%%%%%%%%%%%%%%%%%%%%%%%%%%%%%%%%%%%%%%%%%

\subsubsection{Using the Wrong Time Level in Euler's Method}

Forward Euler uses

\[
\boxed{
f(t_n,y_n).
}
\]

Backward Euler uses

\[
\boxed{
f(t_{n+1},y_{n+1}).
}
\]

The two methods are fundamentally different.

%%%%%%%%%%%%%%%%%%%%%%%%%%%%%%%%%%%%%%%%%%%%%%%%%%%%%%%%%%%%
\subsubsection{Confusing Local and Global Error}

A one-step method may have local defect

\[
O(h^{p+1})
\]

while having global error

\[
O(h^p).
\]

The errors accumulate across time steps.

%%%%%%%%%%%%%%%%%%%%%%%%%%%%%%%%%%%%%%%%%%%%%%%%%%%%%%%%%%%%
\subsubsection{Assuming Smaller Step Size Always Fixes Everything}

Reducing \(h\) improves discretization accuracy only until floating-point,
modeling, and other errors become relevant.

It also increases computational cost.

%%%%%%%%%%%%%%%%%%%%%%%%%%%%%%%%%%%%%%%%%%%%%%%%%%%%%%%%%%%%
\subsubsection{Ignoring Stability}

A method can have a high formal order but still produce a useless solution
if the chosen step lies outside its stability region.

%%%%%%%%%%%%%%%%%%%%%%%%%%%%%%%%%%%%%%%%%%%%%%%%%%%%%%%%%%%%
\subsubsection{Using Explicit Euler on a Stiff Problem With a Large Step}

The exact solution may decay smoothly while Euler oscillates or diverges.

The problem is numerical stability rather than physical instability.

%%%%%%%%%%%%%%%%%%%%%%%%%%%%%%%%%%%%%%%%%%%%%%%%%%%%%%%%%%%%
\subsubsection{Assuming Implicit Means Automatically Accurate}

Implicit methods may be stable at large steps but still have substantial
truncation error.

Stability and accuracy are different properties.

%%%%%%%%%%%%%%%%%%%%%%%%%%%%%%%%%%%%%%%%%%%%%%%%%%%%%%%%%%%%
\subsubsection{Ignoring Nonlinear Solver Error}

An implicit step is only as accurate as the nonlinear equation solve.

Stopping Newton iteration too early introduces additional error.

%%%%%%%%%%%%%%%%%%%%%%%%%%%%%%%%%%%%%%%%%%%%%%%%%%%%%%%%%%%%
\subsubsection{Ignoring Jacobian Quality}

Poor Jacobian approximations can make stiff integration inefficient or
cause nonlinear solves to fail.

%%%%%%%%%%%%%%%%%%%%%%%%%%%%%%%%%%%%%%%%%%%%%%%%%%%%%%%%%%%%
\subsubsection{Using Only Relative Tolerance Near Zero}

When a solution component approaches zero, relative error alone becomes
poorly scaled.

An absolute tolerance is also needed.

%%%%%%%%%%%%%%%%%%%%%%%%%%%%%%%%%%%%%%%%%%%%%%%%%%%%%%%%%%%%
\subsubsection{Using Only Absolute Tolerance for Large Variables}

A single absolute tolerance can demand excessive accuracy for large
variables or insufficient accuracy for small ones.

Relative and absolute tolerances should be combined.

%%%%%%%%%%%%%%%%%%%%%%%%%%%%%%%%%%%%%%%%%%%%%%%%%%%%%%%%%%%%
\subsubsection{Interpolating Events From Coarse Output Only}

An event may occur between stored output times.

Proper event location should use the solver's internal interpolation or
dense-output mechanism.

%%%%%%%%%%%%%%%%%%%%%%%%%%%%%%%%%%%%%%%%%%%%%%%%%%%%%%%%%%%%
\subsubsection{Ignoring Conserved Quantities}

A trajectory can look visually plausible while gradually violating mass,
energy, or population conservation.

Check known invariants.

%%%%%%%%%%%%%%%%%%%%%%%%%%%%%%%%%%%%%%%%%%%%%%%%%%%%%%%%%%%%
\subsubsection{Using Generic Methods for Very Long Hamiltonian Simulations}

Small per-step geometric errors may accumulate systematically.

Structure-preserving methods may provide much better long-time behavior.

%%%%%%%%%%%%%%%%%%%%%%%%%%%%%%%%%%%%%%%%%%%%%%%%%%%%%%%%%%%%
\subsubsection{Applying a Multistep Formula Without Startup Values}

A \(k\)-step method requires sufficient previous solution values.

A separate startup procedure is needed.

%%%%%%%%%%%%%%%%%%%%%%%%%%%%%%%%%%%%%%%%%%%%%%%%%%%%%%%%%%%%
\subsubsection{Ignoring Zero-Stability in Multistep Methods}

A formally accurate multistep formula can still fail to converge if its
homogeneous recurrence is unstable.

%%%%%%%%%%%%%%%%%%%%%%%%%%%%%%%%%%%%%%%%%%%%%%%%%%%%%%%%%%%%
\subsubsection{Treating a Boundary-Value Problem Like an Initial-Value Problem}

Boundary conditions at both endpoints do not usually determine all initial
data directly.

Shooting, finite differences, or collocation may be required.

%%%%%%%%%%%%%%%%%%%%%%%%%%%%%%%%%%%%%%%%%%%%%%%%%%%%%%%%%%%%
\subsection{Applications}
\label{subsec:numerical-ode-applications}
%%%%%%%%%%%%%%%%%%%%%%%%%%%%%%%%%%%%%%%%%%%%%%%%%%%%%%%%%%%%

\begin{applicationbox}

\textbf{Mechanics.}
Numerical ODE methods simulate trajectories, oscillators, rigid-body
motion, and mechanical systems.

\medskip

\textbf{Biology.}
Population, epidemic, ecological, and biochemical models are often systems
of nonlinear ODEs.

\medskip

\textbf{Chemistry.}
Chemical reaction networks can produce strongly stiff ODE systems because
reaction rates span many time scales.

\medskip

\textbf{Engineering.}
Circuit models, control systems, thermal models, and dynamic processes are
integrated numerically.

\medskip

\textbf{Optimal Control.}
Each candidate control trajectory requires integration of state dynamics,
and indirect methods also integrate costate equations.

\medskip

\textbf{Economics.}
Continuous-time growth, capital, adjustment, and dynamic equilibrium models
often require numerical integration.

\medskip

\textbf{Astronomy.}
Orbital and \(N\)-body systems require accurate long-time numerical
integration.

\medskip

\textbf{Scientific Computing.}
ODE integrators frequently form the time-stepping component of larger
simulation software.

\end{applicationbox}

%%%%%%%%%%%%%%%%%%%%%%%%%%%%%%%%%%%%%%%%%%%%%%%%%%%%%%%%%%%%
\subsection{Exercises}
\label{subsec:numerical-ode-exercises}
%%%%%%%%%%%%%%%%%%%%%%%%%%%%%%%%%%%%%%%%%%%%%%%%%%%%%%%%%%%%

\begin{exercise}[1]
Write the standard first-order IVP form.
\end{exercise}

\begin{exercise}[1]
Define a time step.
\end{exercise}

\begin{exercise}[1]
State the forward Euler method.
\end{exercise}

\begin{exercise}[1]
State the backward Euler method.
\end{exercise}

\begin{exercise}[1]
Define local truncation error.
\end{exercise}

\begin{exercise}[1]
Define global error.
\end{exercise}

\begin{exercise}[1]
Define absolute stability.
\end{exercise}

\begin{exercise}[1]
Define stiffness conceptually.
\end{exercise}

\begin{exercise}[1]
Define an adaptive time-step method.
\end{exercise}

\begin{exercise}[1]
Define a boundary-value problem.
\end{exercise}

\begin{exercise}[2]
Use forward Euler to approximate the solution of

\[
y'=2y,
\qquad
y(0)=1
\]

for two time steps with a specified \(h\).
\end{exercise}

\begin{exercise}[2]
Compare the numerical values with

\[
y(t)=e^{2t}.
\]
\end{exercise}

\begin{exercise}[2]
Derive forward Euler from Taylor's theorem.
\end{exercise}

\begin{exercise}[3]
Derive the local truncation error of forward Euler.
\end{exercise}

\begin{exercise}[3]
Perform a convergence study for Euler's method.
\end{exercise}

\begin{exercise}[2]
Apply backward Euler to

\[
y'=-4y.
\]
\end{exercise}

\begin{exercise}[3]
Derive the explicit formula for the backward Euler update in the previous
problem.
\end{exercise}

\begin{exercise}[3]
Compare forward and backward Euler for several step sizes.
\end{exercise}

\begin{exercise}[3]
Apply the trapezoidal method to a linear decay equation.
\end{exercise}

\begin{exercise}[2]
Perform one improved Euler step.
\end{exercise}

\begin{exercise}[2]
Perform one explicit midpoint step.
\end{exercise}

\begin{exercise}[3]
Compare the errors of Euler, midpoint, and improved Euler.
\end{exercise}

\begin{exercise}[2]
Perform one classical RK4 step.
\end{exercise}

\begin{exercise}[3]
Use RK4 for several steps and compare with an exact solution.
\end{exercise}

\begin{exercise}[3]
Write the Butcher tableau for forward Euler.
\end{exercise}

\begin{exercise}[3]
Write the Butcher tableau for the explicit midpoint method.
\end{exercise}

\begin{exercise}[3]
Interpret the stages of RK4.
\end{exercise}

\begin{exercise}[2]
Apply the two-step Adams--Bashforth method when the required history is
given.
\end{exercise}

\begin{exercise}[3]
Apply a predictor--corrector pair.
\end{exercise}

\begin{exercise}[3]
Explain why multistep methods require startup values.
\end{exercise}

\begin{exercise}[3]
Apply the two-step BDF formula to a linear ODE.
\end{exercise}

\begin{exercise}[3]
Construct the first characteristic polynomial of a given linear multistep
method.
\end{exercise}

\begin{exercise}[3]
Check its root condition for zero-stability.
\end{exercise}

\begin{exercise}[2]
Apply forward Euler to

\[
y'=\lambda y
\]

and derive its stability function.
\end{exercise}

\begin{exercise}[3]
Sketch or describe the forward Euler absolute-stability region.
\end{exercise}

\begin{exercise}[3]
Determine the maximum stable step for

\[
y'=-50y.
\]
\end{exercise}

\begin{exercise}[3]
Derive the backward Euler stability function.
\end{exercise}

\begin{exercise}[3]
Show that backward Euler is stable for every negative real \(h\lambda\).
\end{exercise}

\begin{exercise}[3]
Derive the trapezoidal stability function.
\end{exercise}

\begin{exercise}[3]
Explain why the trapezoidal method is A-stable but not L-stable.
\end{exercise}

\begin{exercise}[3]
Explain why stiff problems can force explicit methods to use unnecessarily
small steps.
\end{exercise}

\begin{exercise}[2]
Convert

\[
y''+4y=0
\]

to a first-order system.
\end{exercise}

\begin{exercise}[3]
Apply one Euler step to the resulting system.
\end{exercise}

\begin{exercise}[3]
Write the Jacobian for a given nonlinear ODE system.
\end{exercise}

\begin{exercise}[3]
Construct the Newton system required by backward Euler.
\end{exercise}

\begin{exercise}[3]
Perform one Newton iteration for an implicit Euler step.
\end{exercise}

\begin{exercise}[3]
Explain why Jacobian reuse can reduce stiff-solver cost.
\end{exercise}

\begin{exercise}[2]
Explain the purpose of embedded Runge--Kutta pairs.
\end{exercise}

\begin{exercise}[3]
Given two embedded approximations, compute a local error estimate.
\end{exercise}

\begin{exercise}[3]
Compute a normalized error using specified absolute and relative
tolerances.
\end{exercise}

\begin{exercise}[3]
Determine whether an adaptive step should be accepted.
\end{exercise}

\begin{exercise}[3]
Compute a proposed new step size from an error estimate.
\end{exercise}

\begin{exercise}[3]
Explain why both absolute and relative tolerances are useful.
\end{exercise}

\begin{exercise}[3]
Construct an event function for an object hitting the ground.
\end{exercise}

\begin{exercise}[3]
Explain the difference between terminal and nonterminal events.
\end{exercise}

\begin{exercise}[3]
Check whether Euler preserves the energy of the harmonic oscillator.
\end{exercise}

\begin{exercise}[3]
Explain the motivation for symplectic integrators.
\end{exercise}

\begin{exercise}[3]
Check whether a numerical epidemic solution preserves total population.
\end{exercise}

\begin{exercise}[3]
Formulate a shooting method for a second-order boundary-value problem.
\end{exercise}

\begin{exercise}[3]
Construct the finite-difference equations for

\[
y''=f(x,y)
\]

on a uniform grid.
\end{exercise}

\begin{exercise}[3]
Explain multiple shooting.
\end{exercise}

\begin{exercise}[3]
Explain the method of lines.
\end{exercise}

\begin{exercise}[3]
Estimate observed convergence order from two ODE solution errors.
\end{exercise}

\begin{exercise}[3]
Compare computational cost of Euler and RK4 for a specified number of
steps.
\end{exercise}

\begin{exercise}[4]
Prove first-order convergence of Euler's method under a Lipschitz
condition.
\end{exercise}

\begin{exercise}[4]
Derive a discrete Gronwall error estimate.
\end{exercise}

\begin{exercise}[4]
Derive order conditions for second-order explicit Runge--Kutta methods.
\end{exercise}

\begin{exercise}[4]
Derive the fourth-order classical Runge--Kutta order conditions
conceptually.
\end{exercise}

\begin{exercise}[4]
Study general Runge--Kutta Butcher trees conceptually.
\end{exercise}

\begin{exercise}[4]
Derive Adams--Bashforth formulas by polynomial interpolation of the
derivative.
\end{exercise}

\begin{exercise}[4]
Derive Adams--Moulton formulas.
\end{exercise}

\begin{exercise}[4]
Study consistency conditions for linear multistep methods.
\end{exercise}

\begin{exercise}[4]
Prove the root condition for zero-stability.
\end{exercise}

\begin{exercise}[4]
Study the Dahlquist equivalence theorem.
\end{exercise}

\begin{exercise}[4]
Study the Dahlquist barriers for linear multistep methods.
\end{exercise}

\begin{exercise}[4]
Derive stability functions for several Runge--Kutta methods.
\end{exercise}

\begin{exercise}[4]
Study A-stability and L-stability for implicit methods.
\end{exercise}

\begin{exercise}[4]
Investigate stiffness using Jacobian eigenvalues.
\end{exercise}

\begin{exercise}[4]
Study Rosenbrock methods for stiff ODEs.
\end{exercise}

\begin{exercise}[4]
Study implicit Runge--Kutta methods.
\end{exercise}

\begin{exercise}[4]
Study adaptive embedded Runge--Kutta pairs.
\end{exercise}

\begin{exercise}[4]
Study stiffness-aware variable-order BDF methods.
\end{exercise}

\begin{exercise}[4]
Derive symplectic Euler from Hamiltonian equations.
\end{exercise}

\begin{exercise}[4]
Study the Störmer--Verlet method.
\end{exercise}

\begin{exercise}[4]
Compare long-time energy behavior of RK4 and a symplectic method.
\end{exercise}

\begin{exercise}[4]
Study positivity-preserving time integration.
\end{exercise}

\begin{exercise}[4]
Study numerical differential-algebraic equations.
\end{exercise}

\begin{exercise}[4]
Investigate collocation methods for ODE boundary-value problems.
\end{exercise}

%%%%%%%%%%%%%%%%%%%%%%%%%%%%%%%%%%%%%%%%%%%%%%%%%%%%%%%%%%%%
\subsection{Conceptual Summary}
\label{subsec:numerical-ode-summary}
%%%%%%%%%%%%%%%%%%%%%%%%%%%%%%%%%%%%%%%%%%%%%%%%%%%%%%%%%%%%

Numerical ODE methods approximate solutions of

\[
\boxed{
y'=f(t,y),
\qquad
y(t_0)=y_0.
}
\]

Forward Euler is

\[
\boxed{
y_{n+1}
=
y_n
+
hf(t_n,y_n),
}
\]

while backward Euler is

\[
\boxed{
y_{n+1}
=
y_n
+
hf(t_{n+1},y_{n+1}).
}
\]

Classical RK4 uses four slope evaluations:

\[
\boxed{
y_{n+1}
=
y_n
+
\frac{h}{6}
(
k_1+2k_2+2k_3+k_4
).
}
\]

Numerical ODE methods must be analyzed in terms of

\[
\boxed{
\text{consistency}
+
\text{stability}
+
\text{convergence}.
}
\]

For the test equation

\[
y'=\lambda y,
\]

stability is determined through

\[
\boxed{
y_{n+1}
=
R(h\lambda)y_n.
}
\]

Stiff problems often favor

\[
\boxed{
\text{implicit methods}
}
\]

with strong stability properties.

Adaptive solvers balance accuracy and efficiency by controlling local error
and varying the step size.

Numerical ODE computation therefore combines

\[
\boxed{
\text{differential equations}
+
\text{approximation}
+
\text{stability theory}
+
\text{nonlinear algebra}
+
\text{adaptive computation}.
}
\]

%%%%%%%%%%%%%%%%%%%%%%%%%%%%%%%%%%%%%%%%%%%%%%%%%%%%%%%%%%%%
\subsection{Section Connections}
\label{subsec:numerical-ode-connections}
%%%%%%%%%%%%%%%%%%%%%%%%%%%%%%%%%%%%%%%%%%%%%%%%%%%%%%%%%%%%

\chapterconnection{
Numerical ODEs apply the error, stability, and convergence framework of
Section~\ref{sec:numerical-analysis} to dynamical systems from Chapter~9.
Implicit methods rely heavily on the linear-system techniques of
Section~\ref{sec:numerical-linear-algebra}, while nonlinear implicit solves
use methods related to Section~\ref{sec:numerical-optimization}. Adaptive
Runge--Kutta and stiff integration methods form the computational
foundation of many mathematical simulations and optimal-control
algorithms. The next section, Numerical PDEs, extends these ideas to
partial differential equations by discretizing space and time using
finite-difference, finite-volume, finite-element, spectral, and
method-of-lines approaches.
}

%%%%%%%%%%%%%%%%%%%%%%%%%%%%%%%%%%%%%%%%%%%%%%%%%%%%%%%%%%%%
% SECTION SOURCE TRACKING
%%%%%%%%%%%%%%%%%%%%%%%%%%%%%%%%%%%%%%%%%%%%%%%%%%%%%%%%%%%%

%------------------------------------------------------------
% 11.4 Numerical Ordinary Differential Equations
%------------------------------------------------------------
%
% Source basis:
%   - Standard numerical ODE theory.
%   - Standard Runge--Kutta and linear multistep methods.
%   - Standard stiff integration and adaptive-step techniques.
%
% Used for:
%   - initial-value problems;
%   - numerical time grids;
%   - forward Euler;
%   - local and global truncation error;
%   - consistency and convergence;
%   - backward Euler;
%   - explicit and implicit methods;
%   - trapezoidal method;
%   - improved Euler;
%   - explicit midpoint;
%   - Runge--Kutta methods;
%   - classical RK4;
%   - Butcher tableaux;
%   - one-step methods;
%   - linear multistep methods;
%   - Adams--Bashforth methods;
%   - Adams--Moulton methods;
%   - predictor--corrector methods;
%   - BDF methods;
%   - zero-stability;
%   - Dahlquist equivalence principle;
%   - linear stability test equation;
%   - amplification functions;
%   - absolute stability regions;
%   - A-stability;
%   - L-stability;
%   - stiffness;
%   - systems of ODEs;
%   - higher-order-to-first-order conversion;
%   - Jacobians;
%   - Newton solves for implicit methods;
%   - adaptive time stepping;
%   - embedded Runge--Kutta methods;
%   - absolute and relative tolerances;
%   - rejected steps;
%   - dense output;
%   - event detection;
%   - invariant and conservation issues;
%   - symplectic integration preview;
%   - positivity preservation;
%   - boundary-value problems;
%   - shooting and multiple shooting;
%   - finite-difference BVP methods;
%   - method-of-lines preview;
%   - convergence studies;
%   - computational cost;
%   - applications and exercises.
%
% Notes:
%   - Numerical PDEs are developed next in Section 11.5.
%   - Monte Carlo Methods follow in Section 11.6.
%   - Scientific Computing follows in Section 11.7.
%   - Differential-equation theory appears in Chapter 9.
%
%%%%%%%%%%%%%%%%%%%%%%%%%%%%%%%%%%%%%%%%%%%%%%%%%%%%%%%%%%%%